\section{Related Work}
\label{sec:related}

\para{LLM-as-a-Judge.}
\citet{zheng2023judging} introduced MT-Bench and Chatbot Arena, establishing that strong LLM judges achieve approximately 80\% agreement with human preferences---comparable to human-human agreement.
However, they also documented systematic biases including position preference, verbosity preference, and self-enhancement.
\citet{wang2023large} showed that LLM evaluation rankings can be manipulated by altering response order and proposed calibration strategies.
\citet{stureborg2024large} identified familiarity bias (preference for lower-perplexity text), skewed rating distributions, and anchoring effects, with low inter-sample agreement indicating inconsistency.
Our work extends this line by testing whether LLM judge consensus specifically misaligns with human interestingness judgments.

\para{Self-Preference and Perplexity Bias.}
\citet{panickssery2024llm} demonstrated that LLMs recognize and favor their own generations, with GPT-4 identifying its outputs at 73.5\% accuracy.
Fine-tuning experiments revealed a linear correlation between self-recognition capability and self-preference strength.
\citet{wataoka2024self} deepened this finding by showing that self-preference is driven by \emph{perplexity preference}: LLMs assign higher scores to lower-perplexity text regardless of authorship.
\citet{koo2023cobbler} benchmarked six cognitive biases across 16 LLMs and found an average rank-biased overlap of only 0.44 between human and machine preferences.
These biases motivate our hypothesis---if LLMs prefer familiar text, they may systematically undervalue novel, surprising content.
Unlike prior work that documents these biases in isolation, we directly test whether they aggregate into an inverse relationship between LLM consensus and human interestingness.

\para{LLM Output Diversity and Creativity.}
A growing body of work documents the homogeneity of LLM outputs.
\citet{echoes2024} introduced the Sui Generis score to quantify plot-level uniqueness and found that LLM stories are dramatically less diverse than human writing, with cross-model echo rates showing that different LLMs converge on similar narrative solutions.
\citet{doshi2024generative} found that while generative AI enhances individual creativity, it reduces collective diversity---AI-assisted stories are better but more similar to each other.
\citet{chakrabarty2023art} showed that LLM-generated text scores well on automated creativity metrics but lacks genuine novelty when analyzed by experts.
\citet{gomez2023confederacy} found through blind evaluation of 65 stories from 13 LLMs that human-written stories are often preferred for originality and humor.
Our work builds on these diversity findings: we test whether the rare ``outliers'' that break LLM homogeneity are disproportionately interesting to humans.

\para{Human vs.\ AI Preferences.}
\citet{wu2023style} showed that both human and LLM evaluators prefer longer, grammatically polished text---even when it contains factual errors---suggesting shared ``style over substance'' biases.
\citet{porter2024ai} found that humans rated AI-generated poems higher than human-written poems, possibly because AI poems are more accessible.
\citet{si2024can} reported that LLM-generated research ideas were rated as more novel but less feasible than expert-written ones.
These mixed results motivate our focus on \emph{interestingness} specifically, a dimension where human-LLM alignment has not been directly studied.
